\documentclass[11pt,a4paper]{article}

%% ── Packages ──────────────────────────────────────────────────────────────
\usepackage{mathptmx}
\usepackage[T1]{fontenc}
\usepackage[utf8]{inputenc}
\usepackage[a4paper, top=0.85in, bottom=0.85in, left=1in, right=1in]{geometry}
\usepackage[table]{xcolor}
\usepackage{graphicx}
\usepackage{setspace}
\usepackage{titlesec}
\usepackage{microtype}
\usepackage{enumitem}
\usepackage{booktabs}
\usepackage{longtable}
\usepackage{array}
\usepackage{multirow}
\usepackage{tabularx}
\usepackage{float}
\usepackage{caption}
\usepackage{fancyhdr}
\usepackage{mdframed}
\usepackage{listings}
\usepackage{multicol}
\usepackage{parskip}
\usepackage[hidelinks]{hyperref}
\usepackage{amsmath}

%% ── Colours ───────────────────────────────────────────────────────────────
\definecolor{primaryblue}{RGB}{37,99,235}
\definecolor{darkblue}{RGB}{30,58,138}
\definecolor{lightblue}{RGB}{219,234,254}
\definecolor{accentgreen}{RGB}{16,185,129}
\definecolor{lightgreen}{RGB}{209,250,229}
\definecolor{lightyellow}{RGB}{254,249,195}
\definecolor{lightorange}{RGB}{255,237,213}
\definecolor{accentorange}{RGB}{234,88,12}
\definecolor{lightpurple}{RGB}{237,233,254}
\definecolor{accentpurple}{RGB}{124,58,237}
\definecolor{lightred}{RGB}{254,226,226}
\definecolor{accentred}{RGB}{185,28,28}
\definecolor{tableheadbg}{RGB}{30,58,138}
\definecolor{tablerowbg}{RGB}{239,246,255}
\definecolor{codebg}{RGB}{248,249,250}
\definecolor{bordercolor}{RGB}{203,213,225}
\definecolor{shadecolor}{RGB}{245,245,245}

%% ── Typography ────────────────────────────────────────────────────────────
\setstretch{1.35}
\setlength{\parindent}{0pt}
\setlength{\parskip}{5pt}

%% ── Section headings ──────────────────────────────────────────────────────
\titleformat{\section}
  {\normalfont\fontsize{13}{15}\selectfont\bfseries\color{darkblue}}
  {\thesection}{0.8em}{}[\vspace{-4pt}\rule{\linewidth}{0.8pt}]
\titleformat{\subsection}
  {\normalfont\fontsize{11}{13}\selectfont\bfseries}
  {\thesubsection}{0.8em}{}
\titleformat{\subsubsection}
  {\normalfont\fontsize{11}{13}\selectfont\itshape\bfseries}
  {\thesubsubsection}{0.8em}{}
\titlespacing*{\section}{0pt}{14pt}{4pt}
\titlespacing*{\subsection}{0pt}{10pt}{3pt}
\titlespacing*{\subsubsection}{0pt}{7pt}{2pt}

%% ── Page style ────────────────────────────────────────────────────────────
\pagestyle{fancy}
\fancyhf{}
\fancyhead[L]{\small\textit{\textbf{Brio DBMS} -- Study Notes}}
\fancyhead[R]{\small\textit{CS1211 Exam Prep}}
\fancyfoot[C]{\thepage}
\renewcommand{\headrulewidth}{0.4pt}

%% ── Listings ──────────────────────────────────────────────────────────────
\lstdefinestyle{sql}{
  language=SQL,
  basicstyle=\ttfamily\small,
  keywordstyle=\color{darkblue}\bfseries,
  commentstyle=\color{gray}\itshape,
  stringstyle=\color{accentgreen},
  numbers=none,
  frame=single, framerule=0.4pt,
  rulecolor=\color{bordercolor},
  backgroundcolor=\color{codebg},
  breaklines=true, showstringspaces=false, tabsize=2,
  morekeywords={DELIMITER,TRIGGER,AFTER,BEFORE,EACH,ROW,BEGIN,END,NEW,OLD,
    SIGNAL,SQLSTATE,MESSAGE_TEXT,START,TRANSACTION,COMMIT,ROLLBACK,
    SAVEPOINT,RELEASE,DECLARE,AUTO_INCREMENT,TINYINT,UNSIGNED,
    ENUM,TIMESTAMP,DEFAULT,CURRENT_TIMESTAMP,INDEX,UNIQUE,REPLACE,VIEW},
}
\lstdefinelanguage{JavaScript}{
  keywords={const,let,var,function,async,await,return,if,else,try,catch,import},
  comment=[l]{//}, morecomment=[s]{/*}{*/},
  morestring=[b]", morestring=[b]', morestring=[b]`,
}
\lstdefinestyle{js}{
  language=JavaScript,
  basicstyle=\ttfamily\small,
  keywordstyle=\color{darkblue}\bfseries,
  commentstyle=\color{gray}\itshape,
  stringstyle=\color{accentgreen},
  numbers=none, frame=single, framerule=0.4pt,
  rulecolor=\color{bordercolor},
  backgroundcolor=\color{codebg},
  breaklines=true, showstringspaces=false, tabsize=2,
}
\lstset{style=sql}

%% ── Study-note box environments ───────────────────────────────────────────
% Blue: concept definition
\newmdenv[linecolor=primaryblue, linewidth=1.5pt,
  leftline=true, rightline=false, topline=false, bottomline=false,
  backgroundcolor=lightblue!30,
  innerleftmargin=8pt, innerrightmargin=8pt,
  innertopmargin=5pt, innerbottommargin=5pt,
  skipabove=6pt, skipbelow=4pt]{defbox}

% Green: key point / takeaway
\newmdenv[linecolor=accentgreen, linewidth=1.5pt,
  leftline=true, rightline=false, topline=false, bottomline=false,
  backgroundcolor=lightgreen!30,
  innerleftmargin=8pt, innerrightmargin=8pt,
  innertopmargin=5pt, innerbottommargin=5pt,
  skipabove=6pt, skipbelow=4pt]{keybox}

% Orange: "why?" design decision
\newmdenv[linecolor=accentorange, linewidth=1.5pt,
  leftline=true, rightline=false, topline=false, bottomline=false,
  backgroundcolor=lightorange!50,
  innerleftmargin=8pt, innerrightmargin=8pt,
  innertopmargin=5pt, innerbottommargin=5pt,
  skipabove=6pt, skipbelow=4pt]{whybox}

% Purple: exam Q&A
\newmdenv[linecolor=accentpurple, linewidth=1.5pt,
  leftline=true, rightline=false, topline=false, bottomline=false,
  backgroundcolor=lightpurple!40,
  innerleftmargin=8pt, innerrightmargin=8pt,
  innertopmargin=5pt, innerbottommargin=5pt,
  skipabove=6pt, skipbelow=4pt]{qabox}

% Red: common mistake / warning
\newmdenv[linecolor=accentred, linewidth=1.5pt,
  leftline=true, rightline=false, topline=false, bottomline=false,
  backgroundcolor=lightred!30,
  innerleftmargin=8pt, innerrightmargin=8pt,
  innertopmargin=5pt, innerbottommargin=5pt,
  skipabove=6pt, skipbelow=4pt]{warnbox}

%% ── Helpers ───────────────────────────────────────────────────────────────
\newcommand{\brio}{\textsc{Brio}}
\newcommand{\kw}[1]{\textbf{\texttt{#1}}}
\newcommand{\term}[1]{\textbf{\textit{#1}}}
\newcommand{\T}[1]{\texttt{#1}}          % table/column name

%% ── ToC depth ─────────────────────────────────────────────────────────────
\setcounter{tocdepth}{2}

%% ══════════════════════════════════════════════════════════════════════════
\begin{document}
%% ══════════════════════════════════════════════════════════════════════════

{%
\thispagestyle{empty}
\centering\Large\bfseries\color{darkblue}
Brio DBMS Project\\[4pt]
\normalsize\mdseries\color{black}
Complete Study Notes --- CS1211 Exam Preparation\\[2pt]
{\small Francis Reuben R \textbullet{} Darshan Jain \textbullet{} Harshith BA}\\[2pt]
{\small RV University, Bengaluru \textbullet{} 2025--26}
\par
}

\rule{\linewidth}{1pt}

\tableofcontents

\rule{\linewidth}{0.5pt}

\newpage

%% ══════════════════════════════════════════════════════════════════════════
\section{Project at a Glance}
%% ══════════════════════════════════════════════════════════════════════════

\begin{keybox}
\textbf{\brio} is a full-stack programming learning platform where learners submit code,
earn AI scores, and unlock Kanban project tasks based on lesson mastery. The entire
skill-gate logic --- ``you must pass a lesson to unlock the next task'' --- lives in
the MySQL database as triggers, not in the application code.
\end{keybox}

\subsection{Three-Tier Architecture}

\begin{center}
\renewcommand{\arraystretch}{1.3}
\begin{tabular}{@{}cll@{}}
\toprule
\textbf{Tier} & \textbf{Technology} & \textbf{Role} \\
\midrule
Presentation & Vanilla JS SPA (browser) & Three tabs: Platform / Showcase / Viz \\
Application  & Node.js + Express 5 (\T{server.js}) & 11 REST endpoints, JWT auth, Zod validation \\
Data         & MySQL 8 (\T{brio\_db}) & 13 tables, 3 views, 2 triggers, transactions \\
\bottomrule
\end{tabular}
\end{center}

\subsection{What Makes This Project Non-Trivial}

\begin{itemize}[nosep, leftmargin=1.5em]
  \item \textbf{Skill-gate trigger chain}: two chained triggers fire on every submission INSERT
  \item \textbf{Dual connection pools}: separate read-only MySQL user for the Chaos Sandbox
  \item \textbf{13 tables in 3NF}: strong entities, weak entities, and junction tables
  \item \textbf{15 live SQL queries}: window functions, CTEs, recursive CTEs served via API
  \item \textbf{JWT stateless auth}: token verified on every protected request
\end{itemize}

%% ══════════════════════════════════════════════════════════════════════════
\section{Database Fundamentals (Concepts You Must Know)}
%% ══════════════════════════════════════════════════════════════════════════

\subsection{Relational Model}

\begin{defbox}
A \term{relation} is a table. Each \term{tuple} is a row. Each \term{attribute} is a
column. A \term{schema} is the set of table definitions. A \term{database instance} is
the actual data stored at a point in time.
\end{defbox}

\textbf{Domain constraint}: every value in a column must come from an allowed set (e.g.\
\T{ENUM('learner','admin')} restricts \T{role} to exactly two values).

\subsection{Entity-Relationship (ER) Model}

\begin{defbox}
\textbf{Entity}: a real-world object with independent existence (e.g.\ \T{user}, \T{lesson}).\\
\textbf{Weak entity}: identified only in relation to another entity; its PK includes a
  FK to the owner (e.g.\ \T{kanban\_task} needs \T{proj\_id} to be unique).\\
\textbf{Relationship}: association between entities. \textbf{1:N} (one user owns many
  projects), \textbf{M:N} (many users enroll in many lessons $\to$ junction table).\\
\textbf{Participation}: \textit{total} (every entity participates) vs
  \textit{partial} (some may not).
\end{defbox}

\subsubsection{Brio ER in Plain Language}

\begin{itemize}[nosep, leftmargin=1.5em]
  \item One \T{user} owns many \T{project}s (1:N, partial — not every user has a project)
  \item One \T{project} has many \T{kanban\_task}s (1:N, total — a task must belong to a project)
  \item Many \T{user}s enroll in many \T{lesson}s $\to$ \T{enrollment} junction table (M:N)
  \item One \T{kanban\_task} is gated on one \T{lesson} via \T{lesson\_id} FK
  \item One \T{submission} belongs to one \T{user} + one \T{lesson} (weak entity)
\end{itemize}

\subsection{Keys}

\begin{center}
\renewcommand{\arraystretch}{1.3}
\begin{tabular}{@{}p{3.2cm}p{9.5cm}@{}}
\toprule
\textbf{Key Type} & \textbf{What It Means / Brio Example} \\
\midrule
\textbf{Primary Key (PK)} & Uniquely identifies each row. \T{user\_id INT AUTO\_INCREMENT} in \T{user}. \\
\textbf{Composite PK} & PK spanning multiple columns. \T{(user\_id, lesson\_id)} in \T{enrollment}. \\
\textbf{Foreign Key (FK)} & References the PK of another table. \T{owner\_id} in \T{project} $\to$ \T{user.user\_id}. \\
\textbf{Unique Key} & Not a PK but must be unique. \T{email} and \T{username} in \T{user}. \\
\textbf{Surrogate Key} & Artificial numeric PK (\T{AUTO\_INCREMENT}). Used on strong entities (\T{sub\_id}). \\
\textbf{Natural Key} & Meaningful domain value. Not used in Brio's PKs but \T{email} qualifies. \\
\bottomrule
\end{tabular}
\end{center}

\begin{whybox}
\textbf{Why composite PKs on junction tables?} \T{enrollment(user\_id, lesson\_id)} —
because a \textit{row in enrollment is identified by both values together}. Adding a
surrogate \T{enrollment\_id} would allow duplicate (user, lesson) pairs, which is
semantically wrong. The composite PK enforces the uniqueness constraint for free.
\end{whybox}

\subsection{Normalisation}

The goal: eliminate \textbf{insertion}, \textbf{deletion}, and \textbf{update anomalies}
by ensuring each fact is stored in exactly one place.

\subsubsection{1NF — First Normal Form}
\begin{defbox}
All attribute values are \textbf{atomic} (no lists, no repeating groups). Each column
holds a single indivisible value.\\[4pt]
\textit{Violation example}: storing \T{tags = "python,oop,functions"} in \T{lesson}.\\
\textit{Brio fix}: separate \T{lesson\_tag(lesson\_id, tag)} table — each tag is its
own row.
\end{defbox}

\subsubsection{2NF — Second Normal Form}
\begin{defbox}
1NF + every non-key attribute is \textbf{fully functionally dependent} on the
\textit{entire} primary key (no partial dependency).\\[4pt]
\textit{Only applies to composite PKs.}\\
\textit{Example}: if \T{enrollment} had a column \T{lesson\_title}, that depends only
on \T{lesson\_id} (partial dependency) $\to$ move it to the \T{lesson} table.\\
\textit{Brio}: \T{enrollment(user\_id, lesson\_id, enrolled\_at)} — \T{enrolled\_at}
depends on the full composite PK, so it's in 2NF.
\end{defbox}

\subsubsection{3NF — Third Normal Form}
\begin{defbox}
2NF + no \textbf{transitive dependency} (non-key attribute depends on another non-key
attribute).\\[4pt]
\textit{Example}: if \T{lesson} stored \T{subject\_name} alongside \T{subject\_id},
then \T{lesson\_id $\to$ subject\_id $\to$ subject\_name} is a transitive dependency.\\
\textit{Brio fix}: \T{subject\_name} lives in the \T{subject} table; \T{lesson} only
stores \T{subject\_id} FK. Query it with a JOIN when needed.
\end{defbox}

\begin{keybox}
\textbf{All 13 Brio tables are in 3NF.} The payoff: when a subject name changes, you
update one row in \T{subject} and every lesson that references it automatically reflects
the change.
\end{keybox}

%% ══════════════════════════════════════════════════════════════════════════
\section{The Brio Schema --- All 13 Tables}
%% ══════════════════════════════════════════════════════════════════════════

\begin{center}
\renewcommand{\arraystretch}{1.25}
\begin{tabular}{@{}p{3.5cm}p{2cm}p{7cm}@{}}
\toprule
\rowcolor{tableheadbg}
{\color{white}\textbf{Table}} &
{\color{white}\textbf{Type}} &
{\color{white}\textbf{Purpose \& Key Points}} \\
\midrule
\rowcolor{tablerowbg}
\T{user}              & Strong     & Central entity. PK: \T{user\_id}. \T{role ENUM('learner','admin')}. Passwords bcrypt-hashed. \\
\T{lesson}            & Strong     & Lesson catalogue. FK $\to$ \T{subject}, \T{programming\_language}. \T{difficulty ENUM}. \\
\rowcolor{tablerowbg}
\T{subject}           & Strong     & Subject master (Python, Web Dev\ldots{}). \T{subject\_name UNIQUE}. \\
\T{programming\_language} & Strong & Language master (Python, JS\ldots{}). Referenced by \T{lesson}. \\
\rowcolor{tablerowbg}
\T{project}           & Dependent  & Owned by one user (FK \T{owner\_id}). \T{status ENUM('active','shipped','archived')}. \\
\T{kanban\_task}      & Weak       & Linked to \T{project} + gated on \T{lesson}. \T{status ENUM('locked','in\_progress','done')}. \\
\rowcolor{tablerowbg}
\T{submission}        & Weak       & Event source for triggers. Stores \T{ai\_score TINYINT(0--100)}. \\
\T{lesson\_tag}       & Multivalued & Tags per lesson (1NF compliance). PK: \T{(lesson\_id, tag)}. \\
\rowcolor{tablerowbg}
\T{project\_milestone} & Dependent & Milestones with optional \T{depends\_on} self-FK (used in Q15 recursive CTE). \\
\T{enrollment}        & Junction   & M:N user$\leftrightarrow$lesson. PK: \T{(user\_id, lesson\_id)}. \\
\rowcolor{tablerowbg}
\T{team\_member}      & Junction   & M:N user$\leftrightarrow$project with \T{role} attribute. \\
\T{lesson\_prerequisite} & Junction & M:N lesson$\leftrightarrow$lesson (prerequisite chain). Used in Q9 self-JOIN. \\
\rowcolor{tablerowbg}
\T{user\_badge}       & Junction   & Badges earned by users. PK: \T{(user\_id, badge\_name)}. \\
\bottomrule
\end{tabular}
\end{center}

\subsection{The Skill-Gate Data Model}

The key relationship that powers everything:

\begin{center}
\T{kanban\_task.lesson\_id} $\to$ \T{lesson.lesson\_id}
\end{center}

This FK says: \textit{``this Kanban task is gated behind this lesson.''}
When a user submits a lesson with \T{ai\_score $\geq$ 70}, the trigger looks for all
\T{kanban\_task} rows where \T{lesson\_id} matches AND \T{proj\_id} belongs to the
submitting user's projects, then flips \T{status} from \T{'locked'} to
\T{'in\_progress'}.

\subsection{FK Actions You Must Know}

\begin{center}
\renewcommand{\arraystretch}{1.2}
\begin{tabular}{@{}p{2.8cm}p{9.5cm}@{}}
\toprule
\textbf{Policy} & \textbf{Behaviour / Where Used in Brio} \\
\midrule
\T{CASCADE}    & Delete/update parent $\to$ delete/update all children. \T{project(owner\_id)} cascades on user delete. \\
\T{RESTRICT}   & Prevent parent delete if children exist. \T{lesson(subject\_id)} — can't delete a subject that has lessons. \\
\T{SET NULL}   & Set FK to NULL on parent delete. \T{kanban\_task(lesson\_id)} — task survives if lesson is deleted (just loses its gate). \\
\bottomrule
\end{tabular}
\end{center}

%% ══════════════════════════════════════════════════════════════════════════
\section{DDL --- Data Definition Language}
%% ══════════════════════════════════════════════════════════════════════════

\subsection{Data Types Used in Brio}

\begin{multicols}{2}
\begin{itemize}[nosep, leftmargin=1.5em]
  \item \kw{INT AUTO\_INCREMENT} — surrogate PKs
  \item \kw{VARCHAR(n)} — strings with max length
  \item \kw{TEXT} — unbounded text (bio, content)
  \item \kw{TINYINT UNSIGNED} — 0–255, used for \T{ai\_score} (0–100)
  \item \kw{ENUM('a','b')} — type-safe categorical values
  \item \kw{TIMESTAMP DEFAULT CURRENT\_TIMESTAMP} — auto-records creation time
  \item \kw{DATE} — date only, no time (deadline)
\end{itemize}
\end{multicols}

\subsection{Constraints at a Glance}

\begin{center}
\renewcommand{\arraystretch}{1.2}
\begin{tabular}{@{}p{2.5cm}p{10cm}@{}}
\toprule
\textbf{Constraint} & \textbf{What It Does} \\
\midrule
\T{PRIMARY KEY}    & Unique + NOT NULL identifier. Can be single or composite. \\
\T{NOT NULL}       & Column must always have a value. \\
\T{UNIQUE}         & Value must be unique across all rows (\T{email}, \T{username}). \\
\T{FOREIGN KEY}    & References PK of another table; enforces referential integrity. \\
\T{ENUM}           & Column value must be one of a fixed set. Database rejects invalid values. \\
\T{DEFAULT}        & Value used if none supplied (\T{role DEFAULT 'learner'}). \\
\T{INDEX}          & Speeds up WHERE/JOIN queries on that column. Not a constraint per se. \\
\bottomrule
\end{tabular}
\end{center}

\subsection{Annotated: The Three Most Important Tables}

\begin{lstlisting}[style=sql, caption={\T{user} table --- central auth entity}]
CREATE TABLE user (
    user_id    INT AUTO_INCREMENT PRIMARY KEY,  -- surrogate PK
    username   VARCHAR(50)  NOT NULL UNIQUE,    -- login handle
    email      VARCHAR(100) NOT NULL UNIQUE,    -- contact + alt login
    password   VARCHAR(255) NOT NULL,           -- NEVER plaintext; bcrypt hash
    role       ENUM('learner','admin') DEFAULT 'learner',
    bio        TEXT,                            -- nullable
    created_at TIMESTAMP DEFAULT CURRENT_TIMESTAMP,
    INDEX idx_user_email (email)                -- speeds up login lookup
) ENGINE=InnoDB;                               -- InnoDB = FK support + transactions
\end{lstlisting}

\begin{lstlisting}[style=sql, caption={\T{kanban\_task} --- the skill-gate pivot}]
CREATE TABLE kanban_task (
    task_id    INT AUTO_INCREMENT PRIMARY KEY,
    proj_id    INT NOT NULL,
    lesson_id  INT,            -- NULL = no gate (free task)
    title      VARCHAR(200) NOT NULL,
    status     ENUM('locked','in_progress','done') DEFAULT 'locked',
    priority   INT DEFAULT 5,
    created_at TIMESTAMP DEFAULT CURRENT_TIMESTAMP,
    FOREIGN KEY (proj_id)   REFERENCES project(proj_id)  ON DELETE CASCADE,
    FOREIGN KEY (lesson_id) REFERENCES lesson(lesson_id) ON DELETE SET NULL
    -- SET NULL so deleting a lesson doesn't cascade-delete all tasks
) ENGINE=InnoDB;
\end{lstlisting}

\begin{lstlisting}[style=sql, caption={\T{submission} --- trigger event source}]
CREATE TABLE submission (
    sub_id       INT AUTO_INCREMENT PRIMARY KEY,
    user_id      INT NOT NULL,
    lesson_id    INT NOT NULL,
    ai_score     TINYINT UNSIGNED NOT NULL,  -- CHECK: 0-100 via UNSIGNED + app validation
    submitted_at TIMESTAMP DEFAULT CURRENT_TIMESTAMP,
    FOREIGN KEY (user_id)   REFERENCES user(user_id)     ON DELETE CASCADE,
    FOREIGN KEY (lesson_id) REFERENCES lesson(lesson_id) ON DELETE RESTRICT
    -- RESTRICT: don't delete a lesson that has submission records
) ENGINE=InnoDB;
\end{lstlisting}

\begin{whybox}
\textbf{Why InnoDB?} The \T{MyISAM} engine doesn't support foreign keys or transactions.
\T{InnoDB} (MySQL's default) supports both, making it the only valid choice for a
schema with 13 FK relationships and ACID-compliant transactions.
\end{whybox}

%% ══════════════════════════════════════════════════════════════════════════
\section{DML --- Data Manipulation Language \& Seed Data}
%% ══════════════════════════════════════════════════════════════════════════

\begin{defbox}
\textbf{Seed data strategy}: insert in dependency order — strong entities first, then
dependent/weak entities, then junction tables. You cannot insert a \T{kanban\_task}
before the \T{project} it belongs to exists.
\end{defbox}

\textbf{Insertion order}: \T{subject} $\to$ \T{programming\_language} $\to$ \T{user} $\to$ \T{lesson}
$\to$ \T{project} $\to$ \T{kanban\_task} $\to$ \T{enrollment} $\to$ \T{submission} $\to$ \T{team\_member}

\textbf{bcrypt password hashing}: passwords are never stored as plaintext.
\T{bcrypt.hash(password, 10)} produces a 60-character hash like
\T{\$2b\$10\$...}. During login, \T{bcrypt.compare(input, hash)} returns
true/false. The salt rounds (10) control how slow the hash is — making brute-force
expensive.

%% ══════════════════════════════════════════════════════════════════════════
\section{The 15 Showcase Queries --- Every Concept Explained}
%% ══════════════════════════════════════════════════════════════════════════

Each query teaches one or more SQL concepts. Know \textit{what} the concept is,
\textit{why} it's used, and \textit{what} the query returns.

\subsection{Q1 — GROUP BY + Aggregate Functions}

\begin{defbox}
\T{GROUP BY} collapses all rows with the same value into one summary row. Aggregate
functions (\T{COUNT}, \T{AVG}, \T{MAX}, \T{MIN}, \T{SUM}) operate over each group.
\T{LEFT JOIN} ensures lessons with \textit{zero} submissions still appear (result:
\T{total\_submissions = 0}).
\end{defbox}

\begin{lstlisting}[style=sql]
SELECT l.lesson_id, l.title,
       COUNT(s.sub_id)           AS total_submissions,
       ROUND(AVG(s.ai_score), 2) AS average_score,
       MAX(s.ai_score)           AS highest_score
FROM lesson l
LEFT JOIN submission s ON l.lesson_id = s.lesson_id
GROUP BY l.lesson_id, l.title
ORDER BY total_submissions DESC;
\end{lstlisting}

\begin{keybox}
\textbf{INNER vs LEFT JOIN}: INNER JOIN would omit lessons with no submissions. LEFT JOIN
keeps all lessons and puts NULL for aggregates when no submissions exist. Always ask:
``do I want rows with no matching records?'' $\to$ YES $\to$ LEFT JOIN.
\end{keybox}

\subsection{Q2 — Multi-Table INNER JOIN}

\begin{defbox}
\T{INNER JOIN} returns only rows that have a match in BOTH tables. Join multiple tables
by chaining \T{JOIN ... ON ...} clauses. Alias each table to keep columns readable.
\end{defbox}

Purpose: learner profiles with their enrolment counts — join \T{user}, \T{enrollment},
and \T{lesson} in one query.

\subsection{Q3 — Subquery with NOT IN}

\begin{defbox}
A \term{subquery} is a SELECT inside another SELECT. \T{NOT IN (SELECT ...)} filters
rows that do NOT appear in the inner result set.
\end{defbox}

\begin{lstlisting}[style=sql]
SELECT lesson_id, title FROM lesson
WHERE lesson_id NOT IN (
    SELECT DISTINCT lesson_id FROM submission
);
-- Returns lessons that NO learner has ever submitted
\end{lstlisting}

\begin{warnbox}
\textbf{Pitfall}: If the subquery returns any NULL, \T{NOT IN} behaves unexpectedly
(everything fails). Use \T{NOT EXISTS} when NULLs are possible. Here \T{lesson\_id}
is a FK — it can't be NULL — so \T{NOT IN} is safe.
\end{warnbox}

\subsection{Q4 — Window Function: RANK()}

\begin{defbox}
A \term{window function} computes a value across a ``window'' of rows related to the
current row, without collapsing rows (unlike GROUP BY).\\[4pt]
\T{RANK() OVER (ORDER BY col DESC)} assigns rank 1 to the highest, 2 to next, etc.
Ties get the same rank; the next rank skips (1, 1, 3...).\\
\T{PARTITION BY} restarts ranking per group.
\end{defbox}

\begin{lstlisting}[style=sql]
SELECT u.username,
       ROUND(AVG(s.ai_score), 2) AS avg_score,
       RANK() OVER (ORDER BY AVG(s.ai_score) DESC) AS performance_rank
FROM user u
JOIN submission s ON u.user_id = s.user_id
WHERE u.role = 'learner'
GROUP BY u.user_id, u.username
ORDER BY performance_rank;
\end{lstlisting}

\begin{keybox}
Window functions (\T{RANK}, \T{DENSE\_RANK}, \T{ROW\_NUMBER}, \T{LAG}, \T{LEAD})
require the \T{OVER} clause. They were introduced in SQL:2003 and are available in
MySQL 8+. This is one reason the project specifies \textbf{MySQL 8}.
\end{keybox}

\subsection{Q5 — Common Table Expression (CTE)}

\begin{defbox}
A \term{CTE} (\T{WITH name AS (...)}) defines a named temporary result set valid for
one query. It improves readability and allows referencing the same subquery multiple
times without repeating it.
\end{defbox}

\begin{lstlisting}[style=sql]
WITH high_performers AS (
    SELECT u.user_id, u.username, u.email,
           ROUND(AVG(s.ai_score), 2) AS avg_score
    FROM user u JOIN submission s ON u.user_id = s.user_id
    GROUP BY u.user_id, u.username, u.email
    HAVING AVG(s.ai_score) >= 85        -- filter: only >=85 avg
)
SELECT hp.*, COUNT(e.lesson_id) AS enrolled_lessons
FROM high_performers hp
LEFT JOIN enrollment e ON hp.user_id = e.user_id
GROUP BY hp.user_id, hp.username, hp.email, hp.avg_score
ORDER BY hp.avg_score DESC;
\end{lstlisting}

\subsection{Q6 — CASE Expression}

\begin{defbox}
\T{CASE WHEN ... THEN ... ELSE ... END} is SQL's if-else. It creates a computed column
from conditional logic. Used here to bucket raw scores into labelled performance bands.
\end{defbox}

\begin{lstlisting}[style=sql]
CASE
    WHEN s.ai_score >= 90 THEN 'Excellent'
    WHEN s.ai_score >= 70 THEN 'Good'
    WHEN s.ai_score >= 50 THEN 'Average'
    ELSE 'Needs Work'
END AS performance_band
\end{lstlisting}

\subsection{Q7 — Correlated Subquery}

\begin{defbox}
A \term{correlated subquery} references a column from the outer query — it is
re-executed once per outer row (unlike a regular subquery which runs once).
Used here to find each learner's \textit{latest} submission per lesson.
\end{defbox}

\begin{lstlisting}[style=sql]
WHERE s.submitted_at = (
    SELECT MAX(s2.submitted_at) FROM submission s2
    WHERE s2.user_id = s.user_id AND s2.lesson_id = s.lesson_id
)
-- s2 references s from the outer query -- that's correlation
\end{lstlisting}

\subsection{Q8 — EXISTS}

\begin{defbox}
\T{EXISTS (subquery)} returns TRUE if the subquery returns at least one row.
Unlike \T{IN}, it short-circuits on the first match (faster for large sets).
Used here: ``find learners enrolled in at least one project.''
\end{defbox}

\subsection{Q9 — Self-JOIN}

\begin{defbox}
A \term{self-JOIN} joins a table to itself by aliasing it twice.
Used here on \T{lesson\_prerequisite} to display prerequisite chains:
``Lesson B requires Lesson A'' $\to$ join lesson to itself via the prerequisite table.
\end{defbox}

\subsection{Q10 — Date Functions}

\begin{lstlisting}[style=sql]
WHERE s.submitted_at >= NOW() - INTERVAL 7 DAY
-- Returns submissions from the last 7 days
-- NOW() = current datetime; DATE_SUB(NOW(), INTERVAL 7 DAY) is equivalent
\end{lstlisting}

\subsection{Q11 — LEFT JOIN with Computed Column}

Joins \T{project} to \T{kanban\_task} and computes completion percentage:
\begin{lstlisting}[style=sql]
ROUND(SUM(kt.status = 'done') * 100.0 / NULLIF(COUNT(kt.task_id), 0), 1)
AS completion_pct
-- SUM(condition) counts TRUEs (MySQL treats TRUE=1, FALSE=0)
-- NULLIF prevents division by zero (returns NULL instead)
\end{lstlisting}

\subsection{Q12 — HAVING with Aggregate Condition}

\begin{defbox}
\T{WHERE} filters rows \textit{before} grouping. \T{HAVING} filters \textit{after}
grouping, so it can use aggregate functions.\\
``Find subjects where the average score across all submissions is below 60.''
\end{defbox}

\begin{lstlisting}[style=sql]
GROUP BY l.subject_id, s.subject_name
HAVING AVG(sub.ai_score) < 60
\end{lstlisting}

\subsection{Q13 — UNION}

\begin{defbox}
\T{UNION} stacks two result sets vertically. Both must have the same number of columns
with compatible types. \T{UNION} removes duplicates; \T{UNION ALL} keeps them.
Used here: combine learner and admin rosters into a single unified list.
\end{defbox}

\subsection{Q14 — DENSE\_RANK()}

\begin{defbox}
\T{DENSE\_RANK()} is like \T{RANK()} but with \textit{no gaps}: ties get the same rank
and the next rank is consecutive (1, 1, 2... vs RANK's 1, 1, 3...).
Used here: top 3 lessons by submission volume.
\end{defbox}

\subsection{Q15 — Recursive CTE}

\begin{defbox}
A \term{recursive CTE} calls itself. Structure:
\begin{itemize}[nosep, leftmargin=1.5em]
  \item \textbf{Anchor member}: the base case (starting rows)
  \item \textbf{Recursive member}: joins the CTE to itself to go one level deeper
  \item Terminates when no new rows are added
\end{itemize}
Used here: traverse the \T{project\_milestone} dependency chain — milestone B depends
on milestone A, which depends on milestone Root.
\end{defbox}

\begin{lstlisting}[style=sql]
WITH RECURSIVE milestone_chain AS (
    -- Anchor: root milestones (no dependencies)
    SELECT milestone_id, title, depends_on, 0 AS depth
    FROM project_milestone
    WHERE depends_on IS NULL
    UNION ALL
    -- Recursive: each milestone that depends on one already in the CTE
    SELECT pm.milestone_id, pm.title, pm.depends_on, mc.depth + 1
    FROM project_milestone pm
    JOIN milestone_chain mc ON pm.depends_on = mc.milestone_id
)
SELECT * FROM milestone_chain ORDER BY depth, milestone_id;
\end{lstlisting}

%% ══════════════════════════════════════════════════════════════════════════
\section{Database Views}
%% ══════════════════════════════════════════════════════════════════════════

\begin{defbox}
A \term{view} is a named, stored SELECT query. When you query a view, MySQL executes
the underlying SELECT dynamically — views always reflect current data.
Benefits: encapsulate complex JOINs, simplify API queries, restrict column visibility.
\end{defbox}

\begin{whybox}
\textbf{Why use views in Brio?} Each API endpoint executes one simple \T{SELECT * FROM
view WHERE user\_id = ?}. Without views, every endpoint would repeat complex 4-table
JOINs. Views also make the backend code readable and testable independently of the SQL.
\end{whybox}

\subsection{View 1: learner\_dashboard\_view}

Used by \T{GET /api/profile/me}. Aggregates per-learner stats.

\begin{lstlisting}[style=sql]
CREATE OR REPLACE VIEW learner_dashboard_view AS
SELECT u.user_id, u.username, u.email, u.bio,
    COUNT(DISTINCT e.lesson_id)  AS enrolled_courses,  -- DISTINCT avoids dup count
    ROUND(AVG(s.ai_score), 2)    AS avg_progress,
    COUNT(s.sub_id)              AS total_submissions,
    ROUND(AVG(CASE WHEN s.ai_score IS NOT NULL
                   THEN s.ai_score END), 2) AS avg_score
FROM user u
LEFT JOIN enrollment e ON u.user_id = e.user_id   -- LEFT: include users with 0 enrollments
LEFT JOIN submission  s ON u.user_id = s.user_id  -- LEFT: include users with 0 submissions
WHERE u.role = 'learner'
GROUP BY u.user_id, u.username, u.email, u.bio;
\end{lstlisting}

\subsection{View 2: project\_kanban\_view}

Used by \T{GET /api/projects}. Per-project task breakdown.

\begin{lstlisting}[style=sql]
CREATE OR REPLACE VIEW project_kanban_view AS
SELECT p.proj_id, p.title, p.status, p.deadline,
    COUNT(kt.task_id)                             AS total_tasks,
    SUM(kt.status = 'done')                       AS done_tasks,
    SUM(kt.status = 'in_progress')                AS in_progress_tasks,
    SUM(kt.status = 'locked')                     AS locked_tasks,
    ROUND(SUM(kt.status = 'done') * 100.0
          / NULLIF(COUNT(kt.task_id), 0), 1)      AS completion_pct
FROM project p
LEFT JOIN kanban_task kt ON p.proj_id = kt.proj_id
GROUP BY p.proj_id, p.title, p.status, p.deadline;
-- completion_pct is DERIVED at query time -- never stored, never stale
\end{lstlisting}

\subsection{View 3: skill\_gate\_compliance\_view}

Used by the Showcase module. Shows \textit{why} each locked task is locked.

\begin{lstlisting}[style=sql]
CREATE OR REPLACE VIEW skill_gate_compliance_view AS
SELECT kt.task_id, kt.title AS task_title, kt.status AS task_status,
    l.title AS gate_lesson, u.username,
    MAX(s.ai_score) AS best_score,
    CASE
        WHEN MAX(s.ai_score) >= 70 THEN 'Eligible to unlock'
        WHEN MAX(s.ai_score) IS NULL THEN 'Lesson not submitted yet'
        ELSE 'Score too low (< 70 required)'
    END AS gate_status
FROM kanban_task kt
JOIN project p ON kt.proj_id   = p.proj_id
JOIN user    u ON p.owner_id   = u.user_id
JOIN lesson  l ON kt.lesson_id = l.lesson_id
LEFT JOIN submission s ON s.lesson_id = kt.lesson_id AND s.user_id = u.user_id
WHERE kt.status = 'locked'
GROUP BY kt.task_id, kt.title, kt.status, l.title, u.username;
\end{lstlisting}

%% ══════════════════════════════════════════════════════════════════════════
\section{Triggers}
%% ══════════════════════════════════════════════════════════════════════════

\begin{defbox}
A \term{trigger} is a stored procedure that fires automatically when a specified DML
event (INSERT, UPDATE, DELETE) occurs on a table.\\[4pt]
Syntax: \T{CREATE TRIGGER name \{BEFORE|AFTER\} \{INSERT|UPDATE|DELETE\} ON table FOR
EACH ROW BEGIN ... END}\\[4pt]
\T{NEW.column} = the new row value (INSERT/UPDATE).\\
\T{OLD.column} = the old row value (UPDATE/DELETE).
\end{defbox}

\subsection{The Two-Trigger Skill-Gate Chain}

\textbf{The complete flow when a learner submits a lesson:}

\begin{enumerate}[nosep, leftmargin=2em]
  \item API receives \T{POST /api/submit} with \T{\{lesson\_id, ai\_score\}}
  \item \T{INSERT INTO submission (user\_id, lesson\_id, ai\_score)}
  \item \textbf{Trigger 1 fires} (AFTER INSERT on \T{submission})
  \item If \T{NEW.ai\_score >= 70}: UPDATE all \T{kanban\_task} rows WHERE \T{status='locked'} AND \T{lesson\_id=NEW.lesson\_id} AND task belongs to submitter's projects $\to$ set \T{status='in\_progress'}
  \item \textbf{Trigger 2 fires} (AFTER UPDATE on \T{kanban\_task})
  \item If ALL tasks in the project are \T{status='done'}: UPDATE \T{project} SET \T{status='shipped'}
\end{enumerate}

\begin{lstlisting}[style=sql, caption={Trigger 1: unlock tasks on qualifying score}]
DELIMITER $$
CREATE TRIGGER trg_unlock_task_after_submission
AFTER INSERT ON submission
FOR EACH ROW
BEGIN
    IF NEW.ai_score >= 70 THEN
        UPDATE kanban_task
        SET    status = 'in_progress'
        WHERE  status    = 'locked'
          AND  lesson_id = NEW.lesson_id       -- matches the lesson just submitted
          AND  proj_id IN (
              SELECT proj_id FROM project
              WHERE  owner_id = NEW.user_id   -- only THIS user's projects
          );
    END IF;
END$$
DELIMITER ;
\end{lstlisting}

\begin{lstlisting}[style=sql, caption={Trigger 2: auto-ship project when all tasks done}]
DELIMITER $$
CREATE TRIGGER trg_auto_ship_project
AFTER UPDATE ON kanban_task
FOR EACH ROW
BEGIN
    DECLARE remaining_tasks INT;
    SELECT COUNT(*) INTO remaining_tasks
    FROM   kanban_task
    WHERE  proj_id = NEW.proj_id AND status != 'done';

    IF remaining_tasks = 0 THEN
        UPDATE project SET status = 'shipped'
        WHERE  proj_id = NEW.proj_id;
    END IF;
END$$
DELIMITER ;
\end{lstlisting}

\begin{whybox}
\textbf{Why enforce the skill-gate in the database (trigger) rather than in the Node.js
API code?}
\begin{itemize}[nosep, leftmargin=1.5em]
  \item Application code can be bypassed (direct DB insert, script, admin tool)
  \item A trigger fires on \textit{every} INSERT regardless of what caused it
  \item The rule is a \textit{data integrity rule} — it belongs in the data layer
  \item Separation of concerns: the business rule isn't scattered across endpoints
\end{itemize}
\end{whybox}

\begin{warnbox}
\textbf{MySQL limitation}: a trigger cannot modify the same table it fires on
(prevents infinite loops). Trigger 1 modifies \T{kanban\_task}; Trigger 2 fires on
\T{kanban\_task} and modifies \T{project} — no circular reference.
\end{warnbox}

%% ══════════════════════════════════════════════════════════════════════════
\section{Transactions and ACID}
%% ══════════════════════════════════════════════════════════════════════════

\subsection{ACID Properties}

\begin{center}
\renewcommand{\arraystretch}{1.3}
\begin{tabular}{@{}p{2.4cm}p{10cm}@{}}
\toprule
\textbf{Property} & \textbf{Meaning / Brio Example} \\
\midrule
\textbf{Atomicity}    & Either ALL operations in a transaction commit or NONE do. Enrollment + submission must both succeed or both fail. \\
\textbf{Consistency}  & DB moves from one valid state to another. FK constraints remain satisfied after commit. \\
\textbf{Isolation}    & Concurrent transactions don't see each other's intermediate states. Two users submitting simultaneously don't interfere. \\
\textbf{Durability}   & Once committed, changes survive crashes (InnoDB writes to disk via redo log). \\
\bottomrule
\end{tabular}
\end{center}

\subsection{Transaction Syntax}

\begin{lstlisting}[style=sql]
START TRANSACTION;
SAVEPOINT before_enrollment;          -- snapshot point

INSERT INTO enrollment (user_id, lesson_id, enrolled_at)
VALUES (1, 3, NOW());

-- Check for duplicate
IF (SELECT COUNT(*) FROM enrollment WHERE user_id=1 AND lesson_id=3) > 1 THEN
    ROLLBACK TO SAVEPOINT before_enrollment;   -- undo only to savepoint
    SELECT 'Error: duplicate enrollment' AS result;
ELSE
    RELEASE SAVEPOINT before_enrollment;       -- discard savepoint
    COMMIT;                                    -- persist changes
    SELECT 'Enrolled successfully' AS result;
END IF;
\end{lstlisting}

\begin{keybox}
\textbf{SAVEPOINT} allows partial rollback. \T{ROLLBACK TO SAVEPOINT name} undoes
everything after the savepoint but keeps earlier operations.
\T{ROLLBACK} (no savepoint) undoes the entire transaction.
\end{keybox}

%% ══════════════════════════════════════════════════════════════════════════
\section{Backend Architecture}
%% ══════════════════════════════════════════════════════════════════════════

\subsection{Express 5 Monolith}

Everything lives in \T{website/server.js}. One file = all 11 route handlers + middleware
+ pool setup. This is intentional for educational clarity. In production you'd split
into routes/, controllers/, services/, repositories/.

\subsection{Dual Connection Pools}

\begin{center}
\renewcommand{\arraystretch}{1.3}
\begin{tabular}{@{}p{3.5cm}p{3.5cm}p{5.5cm}@{}}
\toprule
\textbf{Pool} & \textbf{MySQL User} & \textbf{Used By} \\
\midrule
Primary (\T{pool})         & read-write user & All 9 protected API endpoints \\
Read-only (\T{readOnlyPool}) & SELECT-only user & \T{/api/showcase/chaos-sandbox} only \\
\bottomrule
\end{tabular}
\end{center}

\begin{lstlisting}[style=js]
const pool = mysql.createPool({          // 10 connections, read-write
  host: process.env.DB_HOST,
  user: process.env.DB_USER,
  connectionLimit: 10,
  // ...
});

const readOnlyPool = mysql.createPool({ // 5 connections, SELECT-only
  user: process.env.DB_READONLY_USER,   // MySQL user has no INSERT/UPDATE/DELETE grants
  connectionLimit: 5,
});
\end{lstlisting}

\begin{whybox}
\textbf{Why a separate read-only pool for the Chaos Sandbox?}
The Chaos Sandbox lets anyone type arbitrary SQL. Two defence layers:
\begin{enumerate}[nosep, leftmargin=1.5em]
  \item Server-side: reject any query not starting with \T{SELECT} (easily bypassed)
  \item Database-level: the MySQL user itself has no \T{INSERT}/\T{UPDATE}/\T{DELETE} grants
\end{enumerate}
Layer 2 is the real guarantee. Even if a crafty user bypasses the Node.js check via
a subquery trick, the DB refuses to execute DML.
\end{whybox}

\subsection{JWT Authentication}

\begin{defbox}
\term{JSON Web Token (JWT)}: a stateless authentication token. Structure:
\texttt{Header.Payload.Signature} (base64-encoded, dot-separated).
\begin{itemize}[nosep, leftmargin=1.5em]
  \item \textbf{Stateless}: the server doesn't store sessions — the token carries the user info
  \item \textbf{Expiry}: Brio tokens expire after 24 hours
  \item \textbf{Signature}: signed with \T{JWT\_SECRET} — tampered tokens fail verification
\end{itemize}
\end{defbox}

\begin{lstlisting}[style=js]
// Login: verify password -> issue token
const token = jwt.sign(
  { userId: user.user_id, role: user.role },
  process.env.JWT_SECRET,
  { expiresIn: '24h' }
);

// Middleware: verify token on every protected route
function authenticate(req, res, next) {
  const token = req.headers['authorization']?.slice(7); // strip "Bearer "
  try {
    const decoded = jwt.verify(token, process.env.JWT_SECRET);
    req.user = { id: decoded.userId, role: decoded.role };
    next();
  } catch {
    res.status(403).json({ error: 'Invalid or expired token' });
  }
}
\end{lstlisting}

\subsection{Zod Input Validation}

\begin{defbox}
\textbf{Zod} validates request bodies at runtime with a schema. If the body doesn't
match, Zod throws a \T{ZodError} which the global error handler converts to
\T{400 Validation failed \{details\}}.
This prevents malformed data reaching the database layer.
\end{defbox}

\begin{lstlisting}[style=js]
const { ai_score, lesson_id } = z.object({
  ai_score:  z.number().int().min(0).max(100),
  lesson_id: z.number().int().positive(),
}).parse(req.body);
// If parse() throws, the global error handler returns 400 automatically
\end{lstlisting}

\subsection{All 11 API Endpoints}

\begin{center}
\renewcommand{\arraystretch}{1.2}
\begin{tabular}{@{}p{5.5cm}p{1.2cm}p{6.3cm}@{}}
\toprule
\rowcolor{tableheadbg}
{\color{white}\textbf{Route}} &
{\color{white}\textbf{Auth}} &
{\color{white}\textbf{Purpose}} \\
\midrule
\rowcolor{tablerowbg}
\T{POST /api/auth/login}              & None & bcrypt verify $\to$ JWT \\
\T{GET /api/profile/me}               & JWT  & learner\_dashboard\_view \\
\rowcolor{tablerowbg}
\T{POST /api/profile/update}          & JWT  & Update \T{bio} \\
\T{GET /api/projects}                 & JWT  & project\_kanban\_view \\
\rowcolor{tablerowbg}
\T{GET /api/kanban/:projectId}        & JWT  & Tasks by priority \\
\T{GET /api/lessons}                  & JWT  & Enrolled lessons + scores \\
\rowcolor{tablerowbg}
\T{POST /api/submit}                  & JWT  & INSERT submission $\to$ triggers fire \\
\T{GET /api/showcase/stats}           & None & Live user/project/submission counts \\
\rowcolor{tablerowbg}
\T{GET /api/showcase/queries}         & None & Descriptor list for 15 queries \\
\T{POST /api/showcase/query}          & None & Execute query by ID (Q1--Q15, V1--V3) \\
\rowcolor{tablerowbg}
\T{POST /api/showcase/chaos-sandbox}  & None & Arbitrary SELECT via read-only pool \\
\bottomrule
\end{tabular}
\end{center}

%% ══════════════════════════════════════════════════════════════════════════
\section{Frontend Architecture}
%% ══════════════════════════════════════════════════════════════════════════

\begin{defbox}
\textbf{Vanilla JS SPA}: no React, no Vue, no build step. All JS runs directly in the
browser from \T{website/public/}. Three ``tabs'' are really three \T{<div>} sections
shown/hidden based on \T{window.location.hash}.
\end{defbox}

\begin{center}
\renewcommand{\arraystretch}{1.2}
\begin{tabular}{@{}p{2.8cm}p{2.5cm}p{7.5cm}@{}}
\toprule
\textbf{Tab} & \textbf{Auth?} & \textbf{What it shows} \\
\midrule
\textbf{Platform}    & JWT required & Learner dashboard: stats, kanban board, lesson list \\
\textbf{Showcase}    & None          & Live query runner (Q1-Q15), view demos, Chaos Sandbox \\
\textbf{System Viz}  & None          & Static architecture diagrams \\
\bottomrule
\end{tabular}
\end{center}

\textbf{Client-side state in \T{localStorage}:}
\begin{itemize}[nosep, leftmargin=1.5em]
  \item \T{jwt\_token} — JWT string (cleared on logout)
  \item \T{user\_profile} — cached profile data
  \item \T{theme} — dark/light preference
\end{itemize}

\textbf{Hash router}: \T{window.addEventListener('hashchange', renderTab)} — no
server-side routing needed; all navigation is client-side.

%% ══════════════════════════════════════════════════════════════════════════
\section{Key Design Decisions --- The ``Why?'' Questions}
%% ══════════════════════════════════════════════════════════════════════════

These are the questions most likely to come up in a viva. Have a crisp one-sentence
answer and a follow-up explanation ready.

\begin{qabox}
\textbf{Q: Why is the skill-gate enforced in a trigger and not in the API code?}\\
\textbf{A}: Application code can be bypassed; triggers cannot. Any INSERT to
\T{submission} — regardless of source — fires the trigger unconditionally. The
pedagogical rule (score $\geq$ 70 unlocks the task) is a data integrity rule and belongs
in the data layer.
\end{qabox}

\begin{qabox}
\textbf{Q: Why use views instead of writing the JOINs directly in the API endpoints?}\\
\textbf{A}: Views encapsulate complex logic in one place. If the schema changes, you
update the view definition, not every endpoint that uses it. They also act as an
abstraction layer — the API just does \T{SELECT * FROM view WHERE id = ?}.
\end{qabox}

\begin{qabox}
\textbf{Q: Why are composite PKs used on junction tables instead of surrogate keys?}\\
\textbf{A}: A surrogate \T{enrollment\_id} would allow duplicate (user, lesson) pairs —
e.g.\ the same user enrolling in the same lesson twice — which is semantically wrong.
The composite PK \T{(user\_id, lesson\_id)} enforces uniqueness for free, directly
reflecting the relational model.
\end{qabox}

\begin{qabox}
\textbf{Q: What is the Chaos Sandbox and why is it safe?}\\
\textbf{A}: It lets anyone type arbitrary SQL SELECT queries. It's safe via two layers:
(1) server-side check rejects any query not starting with SELECT; (2) the MySQL user
for the read-only pool has no DML privileges — even if you bypass the Node.js check,
the database refuses to execute anything other than SELECT.
\end{qabox}

\begin{qabox}
\textbf{Q: Why InnoDB and not MyISAM?}\\
\textbf{A}: InnoDB supports foreign keys and ACID transactions. MyISAM doesn't. Since
Brio has 13 FK relationships and uses transactions for enrollment, InnoDB is the only
viable choice.
\end{qabox}

\begin{qabox}
\textbf{Q: Explain the normalisation of the lesson table.}\\
\textbf{A}: Storing \T{subject\_name} directly in \T{lesson} would be a transitive
dependency (\T{lesson\_id $\to$ subject\_id $\to$ subject\_name}) — that's a 3NF violation.
Brio stores only \T{subject\_id} FK in \T{lesson} and retrieves the name via JOIN.
If the subject name changes, one row in \T{subject} is updated; all lessons reflect
it automatically.
\end{qabox}

\begin{qabox}
\textbf{Q: Why store \T{lesson\_tag} in a separate table instead of a VARCHAR column?}\\
\textbf{A}: Storing \T{tags = "python,oop,closures"} in one column violates 1NF
(non-atomic values). You can't query ``all lessons tagged python'' without string
parsing. The \T{lesson\_tag(lesson\_id, tag)} table gives each tag its own row — fully
atomic, indexable, and queryable.
\end{qabox}

\begin{qabox}
\textbf{Q: What is JWT and why is it stateless?}\\
\textbf{A}: JWT is a self-contained token with a payload (user ID, role) and a
cryptographic signature. The server verifies the signature using the shared
\T{JWT\_SECRET} — it doesn't need to query a sessions table. This scales well
(no session state to distribute) and the 24-hour expiry limits exposure.
\end{qabox}

\begin{qabox}
\textbf{Q: What does ACID mean? Give an example from Brio.}\\
\textbf{A}: \textbf{A}tomicity (enrollment + submission both commit or neither does),
\textbf{C}onsistency (FK constraints stay satisfied), \textbf{I}solation (concurrent
submissions don't interfere), \textbf{D}urability (committed data survives a crash via
InnoDB's redo log). The enrollment transaction with savepoint-based rollback demonstrates
atomicity and consistency.
\end{qabox}

\begin{qabox}
\textbf{Q: What does LEFT JOIN do that INNER JOIN doesn't?}\\
\textbf{A}: LEFT JOIN returns all rows from the left table, with NULL for columns from
the right table when there's no match. INNER JOIN only returns rows that match in both.
In Brio, \T{learner\_dashboard\_view} uses LEFT JOINs so learners with zero enrollments
or zero submissions still appear in the view (with NULLs for their stats).
\end{qabox}

\begin{qabox}
\textbf{Q: Difference between RANK() and DENSE\_RANK()?}\\
\textbf{A}: If two rows tie for rank 2, \T{RANK()} assigns both rank 2 and skips to
rank 4 (gap). \T{DENSE\_RANK()} assigns both rank 2 and continues with rank 3 (no
gap). Brio uses \T{RANK()} in Q4 (learner performance) and \T{DENSE\_RANK()} in Q14
(top 3 lessons).
\end{qabox}

\begin{qabox}
\textbf{Q: What is a correlated subquery? How is it different from a regular subquery?}\\
\textbf{A}: A regular subquery is evaluated once and its result is reused. A correlated
subquery references a column from the outer query — it is re-evaluated once \textit{per
outer row}. Q7 in Brio uses a correlated subquery to find the latest submission timestamp
per learner per lesson, where the inner query references \T{s.user\_id} and
\T{s.lesson\_id} from the outer query.
\end{qabox}

%% ══════════════════════════════════════════════════════════════════════════
\section{Complete Quick-Reference Summary}
%% ══════════════════════════════════════════════════════════════════════════

\subsection{SQL Concepts Covered in Brio (Cheat Sheet)}

\begin{multicols}{2}
\begin{itemize}[nosep, leftmargin=1.5em, itemsep=1pt]
  \item \textbf{DDL}: CREATE TABLE, constraints, ENUM, INDEX
  \item \textbf{DML}: INSERT (seed data, bcrypt)
  \item \textbf{SELECT}: WHERE, ORDER BY, LIMIT
  \item \textbf{JOINs}: INNER, LEFT, SELF, 4-table
  \item \textbf{Aggregates}: COUNT, AVG, MAX, SUM, GROUP BY
  \item \textbf{HAVING}: filter after grouping
  \item \textbf{Subqueries}: scalar, IN, NOT IN, EXISTS
  \item \textbf{Correlated subquery}: re-evaluated per outer row
  \item \textbf{CASE}: conditional computed column
  \item \textbf{Date functions}: NOW(), INTERVAL, DATE\_SUB()
  \item \textbf{UNION}: stack two result sets
  \item \textbf{Window functions}: RANK(), DENSE\_RANK() OVER
  \item \textbf{CTEs}: WITH clause, non-recursive
  \item \textbf{Recursive CTEs}: WITH RECURSIVE
  \item \textbf{Views}: CREATE OR REPLACE VIEW
  \item \textbf{Triggers}: AFTER INSERT, AFTER UPDATE, NEW/OLD
  \item \textbf{Transactions}: START/COMMIT/ROLLBACK/SAVEPOINT
  \item \textbf{NULLIF}: prevent division-by-zero
  \item \textbf{ROUND()}: decimal rounding
  \item \textbf{FK policies}: CASCADE, RESTRICT, SET NULL
\end{itemize}
\end{multicols}

\subsection{The 15 Queries at a Glance}

\begin{center}
\renewcommand{\arraystretch}{1.1}
\begin{tabular}{@{}cll@{}}
\toprule
\rowcolor{tableheadbg}
{\color{white}\textbf{Q}} &
{\color{white}\textbf{Concept}} &
{\color{white}\textbf{One-line Description}} \\
\midrule
\rowcolor{tablerowbg} Q1  & GROUP BY + AVG/COUNT & Submissions per lesson with avg score \\
Q2  & Multi-JOIN          & Learner profiles with enrolment count \\
\rowcolor{tablerowbg} Q3  & NOT IN subquery     & Lessons never submitted by anyone \\
Q4  & RANK() OVER         & Learner leaderboard by avg score \\
\rowcolor{tablerowbg} Q5  & CTE (WITH)          & High-performers (avg $\geq$ 85) \\
Q6  & CASE expression     & Score banding (Excellent/Good/Average/Needs Work) \\
\rowcolor{tablerowbg} Q7  & Correlated subquery & Latest submission per learner per lesson \\
Q8  & EXISTS              & Learners in at least one project \\
\rowcolor{tablerowbg} Q9  & Self-JOIN           & Lesson prerequisite chains \\
Q10 & Date functions      & Submissions in last 7 days \\
\rowcolor{tablerowbg} Q11 & LEFT JOIN + NULLIF  & Projects with task completion \% \\
Q12 & HAVING              & Subjects with avg score below 60 \\
\rowcolor{tablerowbg} Q13 & UNION               & Combined learner + admin roster \\
Q14 & DENSE\_RANK()       & Top 3 lessons by submission volume \\
\rowcolor{tablerowbg} Q15 & Recursive CTE       & Project milestone dependency traversal \\
\bottomrule
\end{tabular}
\end{center}

\subsection{Things Most Likely to Be Asked in Viva}

\begin{enumerate}[nosep, leftmargin=2em, itemsep=2pt]
  \item \textit{Walk me through what happens when a learner submits a lesson.} \\
        $\to$ POST /api/submit $\to$ Zod validates $\to$ INSERT submission $\to$ Trigger 1 fires $\to$
        if score $\geq$ 70, tasks unlocked $\to$ Trigger 2 fires $\to$ if all tasks done, project ships.

  \item \textit{What is the skill-gate mechanism?} \\
        $\to$ kanban\_task has a lesson\_id FK. Trigger 1 checks this FK to decide which
        tasks to unlock. The database enforces the rule, not the app code.

  \item \textit{Why does your project need two connection pools?} \\
        $\to$ Security. The Chaos Sandbox gets a read-only MySQL user. Even if the Node.js
        check is bypassed, the DB user can only SELECT.

  \item \textit{What is 3NF? Give an example from your project.} \\
        $\to$ No transitive dependencies. subject\_name is in the subject table, not lesson.
        lesson stores only subject\_id FK. Avoids update anomalies.

  \item \textit{Explain the difference between a view and a table.} \\
        $\to$ A table stores data. A view stores a query definition — its data is derived
        at query time. Views are always current; you can't directly INSERT into most views.

  \item \textit{What is ACID? Is it guaranteed in Brio?} \\
        $\to$ Yes, because InnoDB engine is used and transactions wrap multi-step operations.
        The enrollment transaction demonstrates all four properties.

  \item \textit{What's the difference between WHERE and HAVING?} \\
        $\to$ WHERE filters rows before grouping; HAVING filters after. HAVING can use
        aggregate functions; WHERE cannot.

  \item \textit{Why do you use LEFT JOIN in the dashboard view?} \\
        $\to$ To include learners who haven't enrolled in any lessons or haven't submitted
        anything yet. INNER JOIN would hide them.

  \item \textit{Explain the recursive CTE in Q15.} \\
        $\to$ The anchor selects root milestones (depends\_on IS NULL). The recursive part
        joins project\_milestone to the CTE itself to find milestones that depend on
        already-found milestones. Terminates when no new rows match.

  \item \textit{How does bcrypt work?} \\
        $\to$ bcrypt(password, saltRounds) generates a one-way hash. On login,
        bcrypt.compare(input, storedHash) re-hashes the input with the embedded salt
        and compares. Even if the DB is leaked, plaintext can't be recovered.
\end{enumerate}

\vspace{8pt}
\begin{keybox}
\textbf{Golden Rule for the Viva}: for every design decision, always say \textit{why},
not just \textit{what}. ``We used a trigger \textit{because} it fires regardless of which
application path inserts the row'' is the answer; ``we used a trigger'' is not.
\end{keybox}

\end{document}
